\documentclass[final,italian]{memoir}
%%%%%%%%%%%%%%%%%%%%%%%%%%%%%%%%%%%%%%%%%%%%%%%%%%%%%%%%%%%%%%%%%%%%%%%%%%%%%%%
%% Hystrix, the Italian Journal of Mammalogy
%%
%% LaTeX source file to generate an issue cover pages
%%
%% version 1.0
%% created by prea@uninsubria.it 20120525
%%
%%%%%%%%%%%%%%%%%%%%%%%%%%%%%%%%%%%%%%%%%%%%%%%%%%%%%%%%%%%%%%%%%%%%%%%%%%%%%%%
%% This program is free software: you can redistribute it and/or modify
%% it under the terms of the GNU General Public License as published by
%% the Free Software Foundation, either version 3 of the License, or
%% (at your option) any later version.
%%
%% This program is distributed in the hope that it will be useful,
%% but WITHOUT ANY WARRANTY; without even the implied warranty of
%% MERCHANTABILITY or FITNESS FOR A PARTICULAR PURPOSE.  See the
%% GNU General Public License for more details.
%%
%% You should have received a copy of the GNU General Public License
%% along with this program.  If not, see <http://www.gnu.org/licenses/>.
%%%%%%%%%%%%%%%%%%%%%%%%%%%%%%%%%%%%%%%%%%%%%%%%%%%%%%%%%%%%%%%%%%%%%%%%%%%%%%%
%% Notes
%%
%% This LaTeX source automatically generates the cover pages for an Hystrix issue.
%% Cover pages consists of:
%% - a front cover page, which is automatically generated by the macro
%%   \hysfrontcover
%% - a front endpaper (i.e. the verso of the cover page) , automatically generated
%%   by \hysfrontendpaper. Should data appearing on this page be changed, change the
%%   \hysfrontendpaper macro in hystrixmacros.tex
%% - a back endpaper (i.e. the recto of the backcover), automatically  generated
%%   by \hysbackendpaper. Should data appearing on this page be changed, change the
%%   \hysbackendpaper macro in hystrixmacros.tex
%% - a back cover page, that must be typeset by hand, using the hysbackcover
%%   environment provided by hystrixmacros.tex
%%
%% To prepare an issue cover:
%% 1 get a copy of this file and place it into the directory containing all the 
%%   needed files to produce the issue.
%% 2 type in values for the following variables:

\newcommand{\volume}{}      % place here the volume number
\newcommand{\issue}{}        % place here the issue number
\newcommand{\volyear}{2015}   % place here the volume year
\newcommand{\ISSN}{} % place here the journal ISSN
\newcommand{\printedmonth}{settembre}  % place here the printing date month 
\newcommand{\printedyear}{\volyear} % place here the printing date year
\newcommand{\printedpress}{} % place here the printer's name

%% 3 prepare the back cover page (see comments and instruction in the appropriate section, below.
%%
%% 4 compile the LateX source to have a PDF file.

\include{gircmacros}

%% the document begins here

%% we need coloured backgrounds in pages
\usepackage{xcolor}
\usepackage{afterpage}
\usepackage{tabularx}


\begin{document}

%% typeset the front matter here
\pagestyle{empty}
\pagecolor{black}\afterpage{\nopagecolor}

%% Front cover
%%%%%%%%%%%%%%%%%%%%%%%%%%%%%%%%%%%%%%%%%%%%%%%%%%%%%%%%%%%%%%%%%%%%%%%%%%%%%%%%
%% \hysfrontcover     % front cover, "prima di copertina"
\begin{tcolorbox}[size=tight,oversize,
  sharp corners,
  colback=white,
  colframe=white,
  left=56pt,
  right=90pt,
  top=10pt,
  fontupper=\Huge]
\begin{tabular*}{200mm}[t]{lcr}
\includegraphics[height=15mm]{logo_DICAM_web_1.jpg} &
\includegraphics[height=18mm]{logotipo_girc_300dpi_10cm.png} &
\includegraphics[height=18mm]{Logo_ATIt_ritaglio.png} \\
\end{tabular*}
\end{tcolorbox}
\vspace{50mm}
\flushleft
\begin{tabular}{p{110mm}p{40mm}}
   & \multirow{7}{*}{\includegraphics[width=75mm]{bat_moon.png}} \\
   & \\
   & \\
   \sffamily\fontsize{30}{30}\selectfont\textcolor{white}{\textbf{III CONVEGNO ITALIANO}} & \\ [2mm]
   \sffamily\fontsize{30}{30}\selectfont\textcolor{white}{\textbf{SUI CHIROTTERI}} & \\ [3mm]
   \sffamily\huge\textcolor{white}{Trento, 9-11 ottobre 2015} & \\
   & \\
\end{tabular}
\vfill
\begin{tabularx}{180mm}{p{42mm}Xp{.5\textwidth}}
  \selectfont\textcolor{white}{\tiny Con il patrocinio di}\newline
  \includegraphics[width=20mm]{ISPRA CENTRATO bianco.png} &
  &
  \\
  && \\
  
    \selectfont\textcolor{white}{\tiny Sponsored by}\newline
    \includegraphics[width=40mm]{Logo_Pettersson_h_ed2_www2_20.png} &
    &
    \flushright
    {\noindent\sffamily\large\textcolor{white}{published by}}\\
    {\noindent\sffamily\Large\textcolor{white}{Gruppo Italiano Ricerca Chirotteri\\Associazione Teriologica Italiana}} \\
  \e
  
  
\end{tabularx}
%\hfill
%\parbox[t]{.5\textwidth}{
%  \flushright
%  {\noindent\sffamily\large\textcolor{white}{published by}}\\
%  {\noindent\sffamily\Large\textcolor{white}{Gruppo Italiano Ricerca Chirotteri\\Associazione Teriologica %Italiana}}
%}
\clearpage

%% front endpaper, "seconda di copertina"
%%%%%%%%%%%%%%%%%%%%%%%%%%%%%%%%%%%%%%%%%%%%%%%%%%%%%%%%%%%%%%%%%%%%%%%%%%%%%%%%
\hysfrontendpaper

\clearpage

%% here goes the frontmatter
%% next goes the live part of the issue

%% back endpaper, "terza di copertina"
%%%%%%%%%%%%%%%%%%%%%%%%%%%%%%%%%%%%%%%%%%%%%%%%%%%%%%%%%%%%%%%%%%%%%%%%%%%%%%%%
%\hysbackendpaper
\hysheader
\hysfooter
\clearpage
%% uncomment this to have "the Yellow Book"
%\pagecolor{yellow}\afterpage{\nopagecolor}

%%%%%%%%%%%%%%%%%%%%%%%%%%%%%%%%%%%%%%%%%%%%%%%%%%%%%%%%%%%%%%%%%%%%%%%%%%%%%%%
%% back cover section
%%
%% This section must be edited by hand, since it is specific for each issue.
%% The backcover section consists of a series of \paperitem elements, contained 
%% by a \begin and \end standard LaTeX environment.
%% Add papers authors, titles and start pages in the order you need them to
%% appear on the back cover.
%% A \paperitem element structure is straightforward:
%% \paperitem{AUTHORS}{TITLE}{START PAGE}
%% AUTHORS: First name, space, Name initial. No spece between name initials if more than one.
%%          Separate authors with commas.
%% TITLE: just the paper title. properly capitalized. Use \emph{} for emphasis.
%% START PAGE: just the page number that apperas on the paper first page.
%%
%% exemple:
%% \paperitem{Author A., Author B.}{A fake title for an example paper}{32}

\begin{hysbackcover}      % back cover, "quarta di copertina"

Riassunti: Comunicazioni e Poster\\

A cura di Mauro \textsc{Mucedda}, Federica \textsc{Roscioni}, Damiano G. \textsc{Preatoni}

\end{hysbackcover}
\clearpage

\end{document}
%%% Local Variables: 
%%% mode: latex
%%% TeX-master: t
%%% End: 
